% Diagrama de clases UML - Sistema de Gestión de Pedidos
\begin{tikzpicture}[
    class/.style={
        rectangle split, rectangle split parts=3, draw=black, fill=white,
        text=black, rectangle split part align=left, font=\small,
        text width=4.4cm, rounded corners=1pt
    },
    every edge/.style={draw=black, thick},
    labelfont/.style={font=\scriptsize, text=black}
]

\node[class] (cliente) {
    \textbf{Cliente}
    \nodepart{second} cédulaRUC: String [id] \\ nombre: String \\ correo: String
    \nodepart{third} +registrarPedido()
};

\node[class, right=3.8cm of cliente] (pedido) {
    \textbf{Pedido}
    \nodepart{second} idPedido: String [id] \\ fecha: Date \\ total: Decimal
    \nodepart{third} +confirmar() \\ +despachar() \\ +entregar() \\ +anular()
};

\node[class, below=2.6cm of pedido] (lineapedido) {
    \textbf{LíneaPedido}
    \nodepart{second} cantidad: int \\ subtotal: Decimal
    \nodepart{third} +calcularSubtotal()
};

\node[class, below=2.6cm of cliente] (producto) {
    \textbf{Producto}
    \nodepart{second} código: String [id] \\ descripción: String \\ precio: Decimal \\ stock: int
    \nodepart{third} +verificarStock()
};

% Asociación Cliente - Pedido
\draw (cliente) -- (pedido)
    node[labelfont, pos=0.12, above] {1}
    node[labelfont, pos=0.88, above] {0..*};

% Composición Pedido - LineaPedido
\draw[-{Diamond[length=8pt,width=6pt]}] (lineapedido) -- (pedido)
    node[labelfont, pos=0.15, right] {1..*}
    node[labelfont, pos=0.85, right] {1};

% Asociación LineaPedido - Producto
\draw (producto) -- (lineapedido)
    node[labelfont, pos=0.12, above] {1}
    node[labelfont, pos=0.88, above] {0..*};

\end{tikzpicture}
